% LaTeXStruct template: faithfulbook v1
% This deterministic style layer is inserted into an OCR-derived book preamble.
% It deliberately contains no title, author, source-page number, or OCR running head.
\usepackage{geometry}
\geometry{paperwidth=@@PAPER_WIDTH_MM@@mm,paperheight=@@PAPER_HEIGHT_MM@@mm,inner=@@INNER_MARGIN_MM@@mm,outer=@@OUTER_MARGIN_MM@@mm,top=@@TOP_MARGIN_MM@@mm,bottom=@@BOTTOM_MARGIN_MM@@mm,headheight=13.6pt,headsep=4.8mm,footskip=8mm}
\usepackage{xcolor}
\definecolor{LSBookInk}{gray}{0.12}
\definecolor{LSBookMid}{gray}{0.42}
\definecolor{LSBookRule}{gray}{0.68}
\definecolor{LSBookWash}{gray}{0.965}
\usepackage{fancyhdr}
\usepackage{titlesec}
\usepackage{etoc}
\usepackage{etoolbox}
\usepackage{caption}
\usepackage{graphicx}
\usepackage{tcolorbox}
\tcbuselibrary{breakable,skins}
\usepackage{amsthm}
% OCR keeps an explicit printed theorem number as amsthm's optional note so
% LaTeX does not renumber it.  The stock plain style encloses that note in
% parentheses (``Theorem (7.7)''); this plain-compatible head spec keeps the
% source number verbatim and directly after the theorem name instead.
\newtheoremstyle{lsbookplain}%
  {\topsep}{\topsep}% space above/below
  {\itshape}{}% body font / indent
  {\normalfont\bfseries}{.}% head font / punctuation
  {5pt plus 1pt minus 1pt}% space after head
  {\thmname{#1}\thmnumber{ #2}\thmnote{ #3}}
\makeatletter
\let\th@plain\th@lsbookplain
\makeatother
\IfFileExists{microtype.sty}{%
  \usepackage{microtype}
  % XeTeX supports protrusion/tracking, but not font expansion.
  \microtypesetup{protrusion=true,expansion=false,tracking=true}%
}{}

% The class option establishes the nominal size; this fixes the requested
% body leading without changing mathematics or OCR text.
\AtBeginDocument{\fontsize{@@BODY_FONT_PT@@pt}{@@BODY_LEADING_PT@@pt}\selectfont}
\setlength{\parindent}{1.15em}
\setlength{\parskip}{0pt plus 0.15pt}
\setlength{\emergencystretch}{1.5em}
\clubpenalty=9000
\widowpenalty=9000
\displaywidowpenalty=9000
\brokenpenalty=7000
\raggedbottom

% OCR books already carry an authoritative hard break at every source-page
% boundary.  The standard book-class \mainmatter clears once more before it
% resets pagination, which can ship an empty roman-numbered page at a chapter
% boundary.  This wrapper suppresses only that redundant internal clear while
% preserving the class's main-matter state and arabic page reset.
\makeatletter
\newcommand{\LSMainMatter}{%
  \let\LSBookSavedClearDoublePage\cleardoublepage
  \let\cleardoublepage\relax
  \mainmatter
  \let\cleardoublepage\LSBookSavedClearDoublePage
}
% A title/front-matter command may select ``plain`` after document start, so a
% bare \thispagestyle{empty} there is not reliable.  Arm a one-shot flag and
% enforce the empty style at the first actual output page; all later pages use
% the normal faithfulbook styles.
\newif\ifLSBookFirstPageEmpty
\newcommand{\LSFirstPageEmpty}{\global\LSBookFirstPageEmptytrue}
\pretocmd{\@outputpage}{%
  \ifLSBookFirstPageEmpty
    \global\LSBookFirstPageEmptyfalse
    \thispagestyle{empty}%
  \fi
}{}{\PackageWarning{faithfulbook}{Could not install first-page style hook}}
\makeatother

% Generated running heads use only the structured chapter/section tree and
% the newly typeset page number.  OCR page markers remain inert comments.
\pagestyle{fancy}
\fancyhf{}
\fancyhead[LE]{\small\sffamily\color{LSBookMid}\thepage}
\fancyhead[RO]{\small\sffamily\color{LSBookMid}\thepage}
\fancyhead[LO]{\small\sffamily\color{LSBookInk}\nouppercase{\rightmark}}
\fancyhead[RE]{\small\sffamily\color{LSBookInk}\nouppercase{\leftmark}}
\renewcommand{\headrulewidth}{0.35pt}
\renewcommand{\footrulewidth}{0pt}
\renewcommand{\chaptermark}[1]{\markboth{\thechapter\enspace #1}{}}
\renewcommand{\sectionmark}[1]{\markright{\thesection\enspace #1}}
\fancypagestyle{plain}{%
  \fancyhf{}%
  \fancyfoot[C]{\small\sffamily\color{LSBookMid}\thepage}%
  \renewcommand{\headrulewidth}{0pt}%
  \renewcommand{\footrulewidth}{0pt}%
}

% Restrained monochrome hierarchy intended for book-length mathematical text.
\titleformat{\chapter}[display]
  {\normalfont\sffamily\color{LSBookInk}}
  {\filleft\small\bfseries\MakeUppercase{\chaptertitlename}\enspace
   {\fontsize{30}{32}\selectfont\color{LSBookMid}\thechapter}}
  {0.7ex}
  {\titlerule\vspace{1.1ex}\Huge\bfseries\raggedright}
  [\vspace{0.8ex}\titlerule]
\titlespacing*{\chapter}{0pt}{-5mm}{7mm}
\titleformat{\section}
  {\normalfont\Large\sffamily\bfseries\color{LSBookInk}}
  {\thesection}{0.7em}{}
\titleformat{\subsection}
  {\normalfont\large\sffamily\bfseries\color{LSBookInk}}
  {\thesubsection}{0.65em}{}
\titleformat{\subsubsection}
  {\normalfont\normalsize\sffamily\bfseries\color{LSBookInk}}
  {\thesubsubsection}{0.6em}{}
\titlespacing*{\section}{0pt}{2.35ex plus .55ex minus .2ex}{.95ex}
\titlespacing*{\subsection}{0pt}{1.9ex plus .4ex minus .15ex}{.7ex}

% A local, second-pass chapter contents.  The transformer inserts this macro
% only after numbered chapters that actually contain section descendants.
\newcommand{\LSChapterContents}{%
  \begingroup
  \etocsetnexttocdepth{subsection}%
  \etocsettocstyle
    {\par\smallskip\noindent\color{LSBookRule}\rule{\linewidth}{0.35pt}%
     \par\smallskip\small\sffamily\color{LSBookInk}}
    {\par\smallskip\noindent\color{LSBookRule}\rule{\linewidth}{0.35pt}%
     \par\medskip}%
  \localtableofcontents
  \endgroup
}

% The statement boxes are intentionally quiet: a light left rule preserves
% the visual rhythm of traditional theorem typography without rewriting text.
\newcommand{\LSBookStyleStatement}[1]{%
  \ifcsname #1\endcsname
    \tcolorboxenvironment{#1}{%
      enhanced,breakable,blanker,
      borderline west={0.55pt}{0pt}{LSBookRule},
      left=2.2mm,right=0mm,top=0.6ex,bottom=0.6ex,
      before skip=1.0ex plus .25ex,after skip=1.0ex plus .25ex}%
  \fi
}
\AtBeginDocument{%
  \LSBookStyleStatement{theorem}%
  \LSBookStyleStatement{lemma}%
  \LSBookStyleStatement{proposition}%
  \LSBookStyleStatement{corollary}%
  \LSBookStyleStatement{conjecture}%
  \LSBookStyleStatement{claim}%
  \LSBookStyleStatement{fact}%
  \LSBookStyleStatement{definition}%
  \LSBookStyleStatement{remark}%
  \LSBookStyleStatement{observation}%
  \LSBookStyleStatement{example}%
  \LSBookStyleStatement{problem}%
  \LSBookStyleStatement{question}%
  \LSBookStyleStatement{exercise}%
}
\AtBeginEnvironment{proof}{\smallskip}
\AtEndEnvironment{proof}{\smallskip}
\makeatletter
\AtBeginDocument{\@ifundefined{qedsymbol}{}{%
  \renewcommand{\qedsymbol}{\rule{0.62em}{0.62em}}}}
\makeatother

% Algorithms keep the established float environment but receive the same
% thin-rule vocabulary.  Captions use conventional book typography.
\makeatletter
\AtBeginDocument{%
  \@ifundefined{algorithm}{}{%
    \floatstyle{ruled}\restylefloat{algorithm}%
    \floatname{algorithm}{Algorithm}%
  }%
  \@ifundefined{urlstyle}{}{\urlstyle{same}}%
}
\makeatother
\captionsetup{font=small,labelfont=bf,labelsep=period,justification=centering,singlelinecheck=true,skip=5pt}
\AtBeginEnvironment{table}{\captionsetup{position=top}}
\setkeys{Gin}{keepaspectratio}
